\documentclass[aps,pre,superscriptaddress,nofootinbib]{revtex4-2}

\usepackage{fancyhdr}
\usepackage[colorlinks=true,linkcolor=black,citecolor=blue,urlcolor=blue]{hyperref}

% Page style with page numbers
\pagestyle{fancy}
\fancyhf{}
\fancyfoot[C]{\thepage}
\renewcommand{\headrulewidth}{0pt}

% Set paragraph spacing
\setlength{\parskip}{1em}

\begin{document}

\title{Aprendizaje por refuerzo en sistemas estratégicos: aplicación a un agente para Pokémon Showdown}
\author{(Draft)}


\maketitle

% Front matter sections
\section*{Portada}
Portada según el modelo oficial, con los datos de identificación del trabajo.

\newpage
\section*{Dedicatoria}
 (opcional)

\newpage
\section*{Agradecimientos}
 (opcional)

\newpage
\begingroup
\setlength{\parskip}{0pt}
\tableofcontents
\endgroup

\section{Índice de ilustraciones}
 (si procede)

\section{Índice de tablas}
 (si procede)

\newpage
\section{Introducción}
En los últimos años ha habido muchas mejoras en el campo de aprendizaje por refuerzo. Con nuevos diseños de arquitecturas y grafos de recompensa se han superado a jugadores profesionales en juegos como el Go~\cite{silver2016go} o el StarCraft II~\cite{vinyals2019starcraft}. Sin embargo, el campo de aprendizaje por refuerzo es todavía un campo abierto por explorar entornos muy complejos.

En este TFG se propone un nuevo enfoque multiagente aplicado a un entorno particular de Pokémon: las batallas aleatorias individuales de la generación IX disponibles en la plataforma Pokémon Showdown~\cite{pokemonShowdown}. Se trata de un juego de suma cero, con información oculta y un alto grado de estocasticidad. Además, el simulador existente no permite repetir batallas ni comenzar desde un punto intermedio, lo cual constituye un desafío adicional. Este problema representa una considerable dificultad de investigación, por lo que se explorarán diferentes técnicas de aprendizaje multiagente y estrategias de reward shaping.

Para entrenar el modelo se usarán distintas técnicas como Actor-Critic y grafos de recompensa complejos para encontrar un modelo y método de entrenamiento que funciona de manera eficiente.

\section{Objetivos}
En este trabajo se adoptarán los objetivos enunciados como marco de referencia para orientar el diseño, la implementación y la evaluación.

\begin{description}
    \item[H1] \textbf{Representación del estado:} añadir características estratégicamente informativas al espacio de observación mejora la eficiencia de aprendizaje y el rendimiento final de la política frente a una representación más básica del estado.

    \item[H2] \textbf{Self-play:} el entrenamiento mediante self-play produce políticas más robustas y con mayor capacidad de generalización frente a oponentes no vistos que el entrenamiento frente a rivales fijos o heurísticos.

    \item[H3] \textbf{Escalado con el tamaño del equipo:} aumentar el número de equipos incrementa la riqueza estratégica del problema, pero también hace más difícil la optimización por el aumento del espacio de estados, acciones y asignación de crédito; por tanto, el rendimiento no tiene por qué mejorar de forma monótona si la arquitectura y el entrenamiento no escalan adecuadamente.
\end{description}

\section{Antecedentes y contexto}

En 1983 se solucionó el primer problema de aprendizaje por refuerzo con una red neuronal~\cite{barto1983invertedpendulum}. Era el problema del péndulo invertido, un sistema de control clásico. Usar redes neuronales se extendió a otros campos más allá del control, como el aprendizaje de juegos.

En 1993 el programa TD-Gammon de Gerald Tesauro~\cite{tesauro1995gammon} alcanzó un nivel de juego que solo estaba por debajo de los mejores jugadores del momento de backgammon. Backgammon es un juego estocástico (usa dados), de dos agentes, sin información oculta y sin actuación simultánea. Otro factor importante es que el estado inicial siempre es el mismo y no hay tantos estados posibles como en otros juegos.

En 2013 se presentó el programa DQN de DeepMind~\cite{mnih2013dqn}, que fue el primer programa de aprendizaje por refuerzo profundo que superó a los humanos en una variedad de juegos de Atari. Son juegos con observaciones continuas, sin información oculta y sin planificación a largo plazo. Todos estos juegos también daban su puntuación durante cada instante, lo que facilitaba el aprendizaje.

En 2019 el programa AlphaStar de DeepMind~\cite{vinyals2019starcraft} superó a jugadores profesionales en el juego de estrategia en tiempo real StarCraft II. Es un juego con información oculta, actuación simultánea, planificación a largo plazo y un espacio de estados y acciones enorme. Las mayores diferencias con nuestro formato de Pokémon son que Starcraft II es en tiempo real y determinista. Otra diferencia crucial es que no tenemos herramientas para usar aprendizaje supervisado como AlphaStar.

Pokémon es una franquicia de videojuegos que se ha extendido a otros medios como el anime, el manga o los juegos de cartas. Tiene varios juegos competitivos (Unite, GO, TCG), pero en el que se centra este TFG es en las batallas de los juegos principales.

Dentro de los juegos principales, hay distintos formatos de batalla. Cada uno tiene sus propias reglas y restricciones. Los Pokémon y movimientos disponibles, si tienes que elegir un equipo de antemano o no, el número de Pokémon que controlas en el campo, el número de jugadores, etc. Todos estos formatos se caracterizan por tener un elevadísmo número de estados, información oculta, actuación simultánea y un alto grado de estocasticidad. El que se estudia en este TFG es el formato de batallas aleatorias individuales de la generación IX \emph{gen9randombattle}, que se caracteriza por que el simulador genera un equipo aleatorio por jugador. Esto significa que hay investigaciones en distintos formatos.

\begin{sloppypar}
En enero de 2024 presentó Jett Wang su tesis~\cite{wang2024pokemon}. Conseguió ser el octavo mejor jugador del formato \emph{gen4randombattle} entrenando PPO con Self-play para después guiar un Monte Carlo. Lo bueno de usar la cuarta generación es que no es tan inestable como las tres primeras y es la que tiene menos estados posibles (menos Pokémon, movimientos, objetos, etc.) que las generaciones posteriores.
\end{sloppypar}

En marzo de 2025 el agente PokéChamp~\cite{karten2025pokechamp} logra jugar al nivel de los 30\%-10\% mejores jugadores humanos en el formato \emph{gen9ou} usando Minimax LLM. En vez de simular los siguientes turnos de la batalla, usan un modelo de lenguaje y también para decidir la siguiente acción.

En junio de 2025 se publicó la herramienta VGC-Bench~\cite{cameronangliss2025vgcbench}. Es una herramienta para entrenar agentes y compararlos en el formato de batalla de dobles de la generación IX. Consiguieron ganarle al jugador experto Aaron Traylor una vez de cada 5 partidas con una pareja de equipos. La diferencia es que para entrenar su agente usaron aprendizaje supervisado a partir de partidas de jugadores humanos. Otro factor importante que descubrieron es la pérdida de rendimiento de todos sus modelos cuando entrenaban con más equipos.


\section{Resumen de la solución propuesta}
Se indicará la solución aportada para el problema presentado.

\section{Marco teórico o práctico}
Herramientas empleadas: se deberá detallar el marco teórico o práctico en el que se va a desarrollar el trabajo. Si procede, se incluirán las herramientas empleadas.

\section{Planificación y seguimiento}
Se deberá incluir alguna evidencia que muestre tanto la planificación del trabajo, con su distribución de fases y tareas, como su comparación con los datos reales obtenidos tras el desarrollo del trabajo.

\section{Principales aportaciones}
Se deberán destacar las aportaciones importantes del trabajo realizado, teniendo en cuenta los objetivos fijados.

\section{Conclusiones}
Se incluirán todas las conclusiones de tipo técnico y personal.

\section{Vías de trabajo futuro}
Se presentarán posibles ampliaciones y trabajos relacionados por realizar.

\section{Referencias}
\bibliographystyle{unsrtnat}
\bibliography{refs}

\section{Anexos (opcionales)}
Se incluirán otros elementos de interés en el TFG que se consideren necesarios para una mejor comprensión del mismo.

\end{document}